% Options for packages loaded elsewhere
% Options for packages loaded elsewhere
\PassOptionsToPackage{unicode}{hyperref}
\PassOptionsToPackage{hyphens}{url}
\PassOptionsToPackage{dvipsnames,svgnames,x11names}{xcolor}
%
\documentclass[
  spanish,
  11pt,
  letterpaper,
  DIV=11,
  numbers=noendperiod]{scrartcl}
\usepackage{xcolor}
\usepackage[margin=2cm]{geometry}
\usepackage{amsmath,amssymb}
\setcounter{secnumdepth}{-\maxdimen} % remove section numbering
\usepackage{iftex}
\ifPDFTeX
  \usepackage[T1]{fontenc}
  \usepackage[utf8]{inputenc}
  \usepackage{textcomp} % provide euro and other symbols
\else % if luatex or xetex
  \usepackage{unicode-math} % this also loads fontspec
  \defaultfontfeatures{Scale=MatchLowercase}
  \defaultfontfeatures[\rmfamily]{Ligatures=TeX,Scale=1}
\fi
\usepackage{lmodern}
\ifPDFTeX\else
  % xetex/luatex font selection
\fi
% Use upquote if available, for straight quotes in verbatim environments
\IfFileExists{upquote.sty}{\usepackage{upquote}}{}
\IfFileExists{microtype.sty}{% use microtype if available
  \usepackage[]{microtype}
  \UseMicrotypeSet[protrusion]{basicmath} % disable protrusion for tt fonts
}{}
\makeatletter
\@ifundefined{KOMAClassName}{% if non-KOMA class
  \IfFileExists{parskip.sty}{%
    \usepackage{parskip}
  }{% else
    \setlength{\parindent}{0pt}
    \setlength{\parskip}{6pt plus 2pt minus 1pt}}
}{% if KOMA class
  \KOMAoptions{parskip=half}}
\makeatother
% Make \paragraph and \subparagraph free-standing
\makeatletter
\ifx\paragraph\undefined\else
  \let\oldparagraph\paragraph
  \renewcommand{\paragraph}{
    \@ifstar
      \xxxParagraphStar
      \xxxParagraphNoStar
  }
  \newcommand{\xxxParagraphStar}[1]{\oldparagraph*{#1}\mbox{}}
  \newcommand{\xxxParagraphNoStar}[1]{\oldparagraph{#1}\mbox{}}
\fi
\ifx\subparagraph\undefined\else
  \let\oldsubparagraph\subparagraph
  \renewcommand{\subparagraph}{
    \@ifstar
      \xxxSubParagraphStar
      \xxxSubParagraphNoStar
  }
  \newcommand{\xxxSubParagraphStar}[1]{\oldsubparagraph*{#1}\mbox{}}
  \newcommand{\xxxSubParagraphNoStar}[1]{\oldsubparagraph{#1}\mbox{}}
\fi
\makeatother


\usepackage{longtable,booktabs,array}
\newcounter{none} % for unnumbered tables
\usepackage{calc} % for calculating minipage widths
% Correct order of tables after \paragraph or \subparagraph
\usepackage{etoolbox}
\makeatletter
\patchcmd\longtable{\par}{\if@noskipsec\mbox{}\fi\par}{}{}
\makeatother
% Allow footnotes in longtable head/foot
\IfFileExists{footnotehyper.sty}{\usepackage{footnotehyper}}{\usepackage{footnote}}
\makesavenoteenv{longtable}
\usepackage{graphicx}
\makeatletter
\newsavebox\pandoc@box
\newcommand*\pandocbounded[1]{% scales image to fit in text height/width
  \sbox\pandoc@box{#1}%
  \Gscale@div\@tempa{\textheight}{\dimexpr\ht\pandoc@box+\dp\pandoc@box\relax}%
  \Gscale@div\@tempb{\linewidth}{\wd\pandoc@box}%
  \ifdim\@tempb\p@<\@tempa\p@\let\@tempa\@tempb\fi% select the smaller of both
  \ifdim\@tempa\p@<\p@\scalebox{\@tempa}{\usebox\pandoc@box}%
  \else\usebox{\pandoc@box}%
  \fi%
}
% Set default figure placement to htbp
\def\fps@figure{htbp}
\makeatother



\ifLuaTeX
\usepackage[bidi=basic,shorthands=off,spanish,es-nodecimaldot,es-tabla]{babel}
\else
\usepackage[bidi=default,shorthands=off,spanish,es-nodecimaldot,es-tabla]{babel}
\fi
\ifLuaTeX
  \usepackage{selnolig} % disable illegal ligatures
\fi


\setlength{\emergencystretch}{3em} % prevent overfull lines

\providecommand{\tightlist}{%
  \setlength{\itemsep}{0pt}\setlength{\parskip}{0pt}}



 


% preamble.tex – Colegio Portaliano (para pdfLaTeX)

\usepackage[T1]{fontenc}
\usepackage[utf8]{inputenc}
\usepackage{lmodern}
\usepackage{newtxtext}
\usepackage[amsthm]{newtxmath}

% Caracteres Unicode problemáticos
\DeclareUnicodeCharacter{2248}{\ensuremath{\approx}}   % ≈
\DeclareUnicodeCharacter{22EE}{\ensuremath{\vdots}}    % ⋮
\DeclareUnicodeCharacter{22EF}{\ensuremath{\cdots}}    % ⋯
\DeclareUnicodeCharacter{22F1}{\ensuremath{\ddots}}    % ⋱

\AtBeginDocument{%
  \decimalpoint
  \deactivatequoting
}

\everymath{\displaystyle}
\KOMAoption{captions}{tableheading}
\makeatletter
\@ifpackageloaded{caption}{}{\usepackage{caption}}
\AtBeginDocument{%
\ifdefined\contentsname
  \renewcommand*\contentsname{Tabla de contenidos}
\else
  \newcommand\contentsname{Tabla de contenidos}
\fi
\ifdefined\listfigurename
  \renewcommand*\listfigurename{Listado de Figuras}
\else
  \newcommand\listfigurename{Listado de Figuras}
\fi
\ifdefined\listtablename
  \renewcommand*\listtablename{Listado de Tablas}
\else
  \newcommand\listtablename{Listado de Tablas}
\fi
\ifdefined\figurename
  \renewcommand*\figurename{Figura}
\else
  \newcommand\figurename{Figura}
\fi
\ifdefined\tablename
  \renewcommand*\tablename{Tabla}
\else
  \newcommand\tablename{Tabla}
\fi
}
\@ifpackageloaded{float}{}{\usepackage{float}}
\floatstyle{ruled}
\@ifundefined{c@chapter}{\newfloat{codelisting}{h}{lop}}{\newfloat{codelisting}{h}{lop}[chapter]}
\floatname{codelisting}{Listado}
\newcommand*\listoflistings{\listof{codelisting}{Listado de Listados}}
\makeatother
\makeatletter
\makeatother
\makeatletter
\@ifpackageloaded{caption}{}{\usepackage{caption}}
\@ifpackageloaded{subcaption}{}{\usepackage{subcaption}}
\makeatother
\usepackage{bookmark}
\IfFileExists{xurl.sty}{\usepackage{xurl}}{} % add URL line breaks if available
\urlstyle{same}
\hypersetup{
  pdftitle={Carta de autorización para grabación y publicación de defensa oral},
  pdflang={es},
  colorlinks=true,
  linkcolor={blue},
  filecolor={Maroon},
  citecolor={Blue},
  urlcolor={Blue},
  pdfcreator={LaTeX via pandoc}}


\title{Carta de autorización para grabación y publicación de defensa
oral}
\author{}
\date{}
\begin{document}
\maketitle


\thispagestyle{empty}
\begin{center}
    \textbf{\Large Colegio Portaliano}\\[0.3cm]
    \textit{San Felipe}\\[0.3cm]
    \textbf{Departamento de Matemática}\\[0.5cm]
    \hrule
\end{center}

\vspace{0.5cm}

\begin{flushright}
    San Felipe, \hspace{3cm} de \hspace{2cm} de 2026
\end{flushright}

\vspace{0.5cm}

\textbf{Estimado/a Apoderado/a:}

Junto con saludarle, me dirijo a usted para informarle que en el marco
de la asignatura
\textbf{Estadística y Probabilidad Descriptiva e Inferencial (EPDI)},
los estudiantes de 3° y 4° medio participarán en una actividad académica
denominada \textbf{“Defensa Final EPDI”}. Esta actividad consiste en la
presentación oral de un proyecto de investigación estadística frente a
un tribunal compuesto por el profesor y sus compañeros.

Con el objetivo de promover las habilidades de comunicación oral y crear
un repositorio educativo de consulta para futuros estudiantes, se ha
considerado \textbf{grabar en video} las defensas y, con la debida
autorización, publicar una selección de las mejores presentaciones en el
canal institucional de YouTube del colegio o en una lista de
reproducción privada del curso.

La grabación se realizará exclusivamente con fines pedagógicos y no
persigue ningún beneficio comercial. Los estudiantes aparecerán
únicamente en el contexto de su exposición académica, vistiendo el
uniforme oficial del colegio.

Por lo anterior, solicito a usted su \textbf{consentimiento explícito}
para que su pupilo/a pueda ser grabado/a y, eventualmente, el video
pueda ser publicado en YouTube. La participación en la grabación es
\textbf{totalmente voluntaria} y no afectará la calificación del
estudiante en caso de no autorizarlo.

Agradezco desde ya su comprensión y colaboración con esta iniciativa que
busca enriquecer el proceso formativo de nuestros estudiantes.

Atentamente,

\rule{5cm}{0.4pt}

\textbf{Prof. Hans Sigrist} Profesor EPDI Colegio Portaliano

\vfill

\begin{center}
    \hrule
    \vspace{0.3cm}
    \textbf{\Large AUTORIZACIÓN}\\
    \textit{(Completar y entregar al profesor)}
    \vspace{0.3cm}
    \hrule
\end{center}

\vspace{0.3cm}

Yo, \underline{\hspace{6cm}}, RUT \underline{\hspace{4cm}}, apoderado/a
de: \newline \underline{\hspace{8cm}} (nombre completo del estudiante),
del curso \underline{\hspace{1cm}}, \newline \underline{\hspace{1cm}}
\textbf{AUTORIZO} \hspace{1cm} \underline{\hspace{1cm}}
\textbf{NO AUTORIZO} (Marque con una X) que mi pupilo/a sea grabado/a en
video durante la actividad de \textbf{Defensa Final EPDI} y que dicho
material pueda ser publicado en la plataforma YouTube con fines
exclusivamente educativos.

\vspace{0.5cm}

Firma del Apoderado/a: \underline{\hspace{6cm}}

\vspace{0.5cm}

Fecha: \underline{\hspace{1cm}} de \underline{\hspace{3cm}} de 2026.

\vfill

\begin{center}
    \textit{Este documento debe ser entregado firmado al profesor de la asignatura antes del día de la defensa.}
\end{center}




\end{document}
